% answers.tex -- marking criteria and sample answers for questions.tex (Module 5 short answer, Set 2).
% NESA marking-guideline style: a criteria table for each part, then a sample answer.
% Every number: checks/m5-short-answer-2/numbers.py
% Build: TEXINPUTS=.: xelatex answers.tex (twice)
\documentclass[12pt]{article}
\usepackage{quiz}
\hypersetup{pdftitle={Short Answer Questions Set 2 - Module 5 - Advanced Mechanics - Marking guidelines}}
% links between the two PDF files are not possible, so question headings are plain text here
\renewcommand\hyperlink[2]{#2}
\newcommand\anstopic[1]{\par\Needspace{8\baselineskip}\vspace{4mm}%
  \begin{center}{\headfont\bfseries\fontsize{14}{18}\selectfont #1\par}\end{center}\vspace{2pt}}

\begin{document}
\quiztitle{Year 12 Physics}{Short Answer Questions (Set 2) - Module 5 - Advanced Mechanics:\\
  Marking Guidelines and Sample Answers}

Data used: $g = \qty{9.8}{m.s^{-2}}$, $G = \qty{6.67e-11}{N.m^2.kg^{-2}}$,
$M_\mathrm{E} = \qty{6.0e24}{kg}$, $r_\mathrm{E} = \qty{6.371e6}{m}$. In calculations, a correct
answer from a correct method with an earlier error carried forward still earns the later marks.

% =====================================================================
\anstopic{Topic 1 - Projectile Motion}

% ---- Q1
\ansq
\anspart{a}
\begin{crit}
  \critrow{Correctly calculates the maximum height above the floor (including the \qty{2.0}{m})}{2}
  \critrow{Calculates the vertical component of the velocity OR the rise above the release point}{1}
\end{crit}
\sample
$u_y = 8.5\sin 52^\circ = \qty{6.70}{m.s^{-1}}$. At the top $v_y = 0$, so the rise is
$s = \dfrac{u_y^2}{2g} = \dfrac{6.70^2}{2 \times 9.8} = \qty{2.29}{m}$.
Maximum height $= 2.0 + 2.29 = \qty{4.3}{m}$ above the floor.
\anspart{b}
\begin{crit}
  \critrow{Correctly calculates the speed}{1}
\end{crit}
\sample
At the top only the horizontal velocity is left: $v = 8.5\cos 52^\circ = \qty{5.2}{m.s^{-1}}$.
\anspart{c}
\begin{crit}
  \critrow{Explains that gravity still acts, so the acceleration is $g$ downwards (only the vertical
    velocity is zero)}{1}
\end{crit}
\sample
Gravity still acts on the ball at the top, so its acceleration is still \qty{9.8}{m.s^{-2}} downwards.
Only the vertical \emph{velocity} is zero at that instant.

% ---- Q2
\ansq
\anspart{a}
\begin{crit}
  \critrow{Derives the time of flight AND uses $R = u_x t$ AND simplifies to $\frac{u^2\sin 2\theta}{g}$}{3}
  \critrow{Derives the time of flight AND writes $R = u\cos\theta \times t$}{2}
  \critrow{Resolves the velocity into components OR finds the time of flight}{1}
\end{crit}
\sample
$u_x = u\cos\theta$ and $u_y = u\sin\theta$. The projectile lands when the vertical displacement is
zero:
\[ 0 = (u\sin\theta)t - \tfrac{1}{2}gt^2 \quad\Rightarrow\quad t = \frac{2u\sin\theta}{g} \]
\[ R = u_x t = u\cos\theta \times \frac{2u\sin\theta}{g} = \frac{u^2(2\sin\theta\cos\theta)}{g}
     = \frac{u^2\sin 2\theta}{g} \]
\anspart{b}
\begin{crit}
  \critrow{Shows that $\sin 2\theta$ is the same for both angles AND links this to equal ranges (or
    explains the trade-off between $u_x$ and time)}{2}
  \critrow{Provides some relevant information}{1}
\end{crit}
\sample
For \ang{30}: $\sin 60^\circ = 0.866$. For \ang{60}: $\sin 120^\circ = 0.866$. With the same $u$ and
$g$, the range equation gives the same $R$. Physically, the \ang{30} launch has a larger horizontal
velocity but a shorter time of flight; the \ang{60} launch has a smaller horizontal velocity but a
longer time of flight. (Any pair of angles that add to \ang{90} gives the same range.)

% ---- Q3
\ansq
\begin{crit}
  \critrow{Describes at least three correct differences, with reference to air resistance opposing
    the motion}{3}
  \critrow{Describes two correct differences}{2}
  \critrow{Describes one correct difference}{1}
\end{crit}
\sample
Air resistance acts opposite to the velocity, so it slows the projectile down.
\begin{itemize}
  \item The horizontal velocity is no longer constant: it decreases during the flight.
  \item The maximum height is lower, and the range is shorter.
  \item The path is no longer a symmetrical parabola: the projectile comes down more steeply than it
        went up, and takes longer to come down than to go up.
\end{itemize}

% ---- Q4
\ansq
\anspart{a}
\begin{crit}
  \critrow{Correctly calculates the horizontal distance}{3}
  \critrow{Correctly calculates the time of fall}{2}
  \critrow{Uses a correct vertical equation with an initial vertical velocity of zero}{1}
\end{crit}
\sample
The package starts with zero vertical velocity: $500 = \frac{1}{2}(9.8)t^2$, so
$t = \sqrt{\dfrac{1000}{9.8}} = \qty{10.1}{s}$. It keeps the plane's horizontal velocity, so the
horizontal distance $= 70 \times 10.1 = \qty{707}{m}$ (about \qty{7.1e2}{m}).
\anspart{b}
\begin{crit}
  \critrow{Correctly describes the path for BOTH observers, with a reason}{2}
  \critrow{Correctly describes the path for one observer}{1}
\end{crit}
\sample
The package and the plane have the same horizontal velocity (\qty{70}{m.s^{-1}}). So the pilot sees
the package fall in a straight line vertically downwards, staying directly below the plane. The
person on the ground sees the package move forwards at a constant \qty{70}{m.s^{-1}} while it speeds
up downwards, so its path is a parabola.

% ---- Q5
\ansq
\anspart{a}
\begin{crit}
  \critrow{Correctly calculates the predicted range}{1}
\end{crit}
\sample
$R = \dfrac{5.0^2 \times \sin 90^\circ}{9.8} = \qty{2.55}{m}$.
\anspart{b}
\begin{crit}
  \critrow{Identifies \ang{45} AND explains that $\sin 2\theta$ has its maximum value of 1 when
    $2\theta = 90^\circ$}{2}
  \critrow{Identifies \ang{45}}{1}
\end{crit}
\sample
\ang{45}. The range is $R = \dfrac{u^2\sin 2\theta}{g}$. With $u$ and $g$ fixed, $R$ is largest when
$\sin 2\theta$ is largest. Its maximum value is 1, when $2\theta = 90^\circ$, so $\theta = 45^\circ$.
The measured data also show the largest range at \ang{45}.
\anspart{c}
\begin{crit}
  \critrow{Correctly estimates the launch speed AND explains why a lower speed reduces every range}{2}
  \critrow{Estimates the launch speed OR explains the effect of a lower speed}{1}
\end{crit}
\sample
At \ang{45}, $\sin 2\theta = 1$, so $R = \dfrac{u^2}{g}$ and $u = \sqrt{Rg} = \sqrt{2.44 \times 9.8} =
\qty{4.9}{m.s^{-1}}$. The range is proportional to $u^2$ at every angle, so a launch speed lower than
\qty{5.0}{m.s^{-1}} makes every range smaller than predicted. This is a systematic error: it affects all
the results in the same direction.

% ---- Q6
\ansq
\anspart{a}
\begin{crit}
  \critrow{Correctly calculates the height}{3}
  \critrow{Calculates the time to reach the wall AND uses a correct vertical equation}{2}
  \critrow{Resolves the velocity into components OR calculates the time}{1}
\end{crit}
\sample
$u_x = 20\cos 60^\circ = \qty{10.0}{m.s^{-1}}$ and $u_y = 20\sin 60^\circ = \qty{17.3}{m.s^{-1}}$.
Time to reach the wall: $t = \dfrac{24}{10.0} = \qty{2.4}{s}$.
\[ s_y = u_y t - \tfrac{1}{2}gt^2 = 17.32 \times 2.4 - 4.9 \times 2.4^2 = 41.57 - 28.22 = \qty{13.3}{m} \]
\anspart{b}
\begin{crit}
  \critrow{States downwards, with a supporting calculation or reason}{1}
\end{crit}
\sample
Downwards. $v_y = u_y - gt = 17.32 - 9.8 \times 2.4 = \qty{-6.2}{m.s^{-1}}$. The vertical velocity is
negative, so the ball is already moving down (it reached its highest point at $t = \frac{17.32}{9.8} =
\qty{1.77}{s}$).

% =====================================================================
\anstopic{Topic 2 - Circular Motion}

% ---- Q7
\ansq
\begin{crit}
  \critrow{Explains that the passenger tends to keep moving in a straight line (inertia), that the car
    turns towards the centre, AND that the door/seat pushes the passenger inwards (there is no
    outward force)}{3}
  \critrow{Gives two of these points}{2}
  \critrow{Gives one relevant point}{1}
\end{crit}
\sample
There is no outward force on the passenger. Because of inertia (Newton's first law), the passenger
tends to keep moving in a straight line. The car turns towards the centre of the curve, so the door
moves into the passenger's path and pushes on them. This inward push (a normal force, with friction
from the seat) provides the centripetal force that makes the passenger turn with the car. The
passenger feels this push from the door, and it feels as if they are being thrown outwards.

% ---- Q8
\ansq
\anspart{a}
\begin{crit}
  \critrow{Correctly calculates the normal force}{2}
  \critrow{Writes $mg - N = \frac{mv^2}{r}$ (net force downwards, towards the centre)}{1}
\end{crit}
\sample
At the top, the centre of the circle is below the car, so the net force is downwards:
$mg - N = \dfrac{mv^2}{r}$.
\[ N = 1500 \times 9.8 - \frac{1500 \times 15^2}{40} = 14\,700 - 8440 = \qty{6.3e3}{N} \text{ (upwards)} \]
\anspart{b}
\begin{crit}
  \critrow{Correctly calculates the maximum speed}{2}
  \critrow{States that contact is lost when $N = 0$, so $mg = \frac{mv^2}{r}$}{1}
\end{crit}
\sample
The car loses contact when the normal force becomes zero. Then the weight alone gives the
centripetal force: $mg = \dfrac{mv^2}{r}$, so $v = \sqrt{gr} = \sqrt{9.8 \times 40} = \qty{19.8}{m.s^{-1}}$.

% ---- Q9
\ansq
\anspart{a}
\begin{crit}
  \critrow{Correctly calculates the torque}{1}
\end{crit}
\sample
$\tau = rF = 1.5 \times (30 \times 9.8) = \qty{441}{N.m}$ (the weight is at \ang{90} to the seesaw).
\anspart{b}
\begin{crit}
  \critrow{Correctly calculates the distance}{2}
  \critrow{Equates the two torques}{1}
\end{crit}
\sample
For balance, the two torques about the pivot are equal and opposite:
$45 \times 9.8 \times r = 441$, so $r = \qty{1.0}{m}$.
\anspart{c}
\begin{crit}
  \critrow{Explains that the torque of the \qty{45}{kg} child increases with distance, so the torques are
    unbalanced and there is a net torque, which causes rotation}{2}
  \critrow{Provides some relevant information}{1}
\end{crit}
\sample
Torque $= rF$. The weight of the \qty{45}{kg} child stays the same, but $r$ increases, so this child's
torque becomes larger than \qty{441}{N.m}. The torques no longer balance, so there is a net torque
about the pivot. A net torque causes the seesaw to rotate: the \qty{45}{kg} child's side goes down.

% ---- Q10
\ansq
\anspart{a}
\begin{crit}
  \critrow{Correctly calculates the angular velocity}{1}
\end{crit}
\sample
$f = \dfrac{1200}{60} = \qty{20}{Hz}$, so $\omega = 2\pi f = \qty{126}{rad.s^{-1}}$.
\anspart{b}
\begin{crit}
  \critrow{Correctly calculates the centripetal acceleration}{2}
  \critrow{Uses $a = \omega^2 r$ or $a = \frac{v^2}{r}$ with a correct value}{1}
\end{crit}
\sample
$a = \omega^2 r = 125.7^2 \times 0.25 = \qty{3.95e3}{m.s^{-2}}$ (about 400 times $g$; \num{3.9e3} or
\num{4.0e3} accepted).
\anspart{c}
\begin{crit}
  \critrow{Explains that the wall provides the centripetal force for the clothes but not for the water,
    so the water moves off in a straight line through the holes}{1}
\end{crit}
\sample
The drum wall pushes the clothes towards the centre, so the clothes move in a circle. Where there is
a hole, nothing pushes the water inwards, so it keeps moving in a straight line (along the tangent),
out through the holes and away from the clothes.

% ---- Q11
\ansq
\anspart{a}
\begin{crit}
  \critrow{Correctly calculates the speed}{1}
\end{crit}
\sample
$v = \dfrac{2\pi r}{T} = \dfrac{2\pi \times 0.80}{0.40} = \qty{12.6}{m.s^{-1}}$.
\anspart{b}
\begin{crit}
  \critrow{Correctly calculates the tension}{2}
  \critrow{Uses $F = \frac{mv^2}{r}$}{1}
\end{crit}
\sample
$F_\mathrm{T} = \dfrac{mv^2}{r} = \dfrac{0.50 \times 12.566^2}{0.80} = \qty{98.7}{N}$ (about \qty{99}{N}).
\anspart{c}
\begin{crit}
  \critrow{Correctly calculates the shortest period}{2}
  \critrow{Uses $F = \frac{4\pi^2 mr}{T^2}$ with $F = \qty{150}{N}$ (or finds the maximum speed)}{1}
\end{crit}
\sample
$F = \dfrac{4\pi^2 mr}{T^2}$, so $T = 2\pi\sqrt{\dfrac{mr}{F}} = 2\pi\sqrt{\dfrac{0.50 \times 0.80}{150}}
= \qty{0.32}{s}$. (Or: $v_\mathrm{max} = \sqrt{\frac{Fr}{m}} = \qty{15.5}{m.s^{-1}}$, then
$T = \frac{2\pi r}{v} = \qty{0.32}{s}$.)
\anspart{d}
\begin{crit}
  \critrow{States that the tension doubles, with a reason ($F = \frac{4\pi^2 mr}{T^2}$, so $F \propto r$
    when $T$ is constant)}{1}
\end{crit}
\sample
The tension doubles. $F = \dfrac{4\pi^2 mr}{T^2}$, so with the same mass and period, $F \propto r$.
(Doubling $r$ also doubles the speed, and $\frac{mv^2}{r}$ doubles.) The new tension would be about
\qty{197}{N}, more than \qty{150}{N}, so the string would break.

% ---- Q12
\ansq
\anspart{a}
\begin{crit}
  \critrow{Correctly calculates all five values of $\frac{1}{r}$}{1}
\end{crit}
\sample
$\frac{1}{r}$ (\si{m^{-1}}): 2.5, 2.0, 1.67, 1.25, 1.0.
\anspart{b}
\begin{crit}
  \critrow{Plots all five points correctly AND draws an appropriate straight line of best fit}{3}
  \critrow{Plots at least four points correctly AND draws a line of best fit}{2}
  \critrow{Plots some points correctly}{1}
\end{crit}
\sample
\begin{center}
  \begin{tikzpicture}
    \begin{axis}[width=105mm, height=72mm, axis lines=left, xmin=0, xmax=3.0, ymin=0, ymax=1.4,
        xtick={0,0.5,...,3.0}, ytick={0,0.2,...,1.4}, grid=major, major grid style={black!25},
        tick label style={font=\small}, label style={font=\small},
        xticklabel style={/pgf/number format/fixed, /pgf/number format/fixed zerofill, /pgf/number format/precision=1},
        yticklabel style={/pgf/number format/fixed, /pgf/number format/fixed zerofill, /pgf/number format/precision=1},
        xlabel={$\frac{1}{r}$ (\si{m^{-1}})}, ylabel={Force $F$ (\si{N})},
        scaled ticks=false, enlargelimits=false]
      \addplot[only marks, mark=x, mark size=3pt, thick] coordinates {(2.5,1.14) (2.0,0.89) (1.667,0.76) (1.25,0.55) (1.0,0.46)};
      \addplot[line width=0.8pt, domain=0:2.9] {0.45*x};
    \end{axis}
  \end{tikzpicture}
\end{center}
\anspart{c}
\begin{crit}
  \critrow{Calculates the gradient, uses gradient $= mv^2$ to find the mass AND compares it with
    \qty{0.050}{kg}}{2}
  \critrow{Calculates the gradient OR states that the gradient equals $mv^2$}{1}
\end{crit}
\sample
Gradient $= \dfrac{1.13 - 0}{2.5 - 0} \approx \qty{0.45}{N.m}$ (accept 0.44 to 0.47). Since
$F = mv^2 \times \dfrac{1}{r}$, the gradient is $mv^2$:
$m = \dfrac{0.45}{3.0^2} = \qty{0.050}{kg}$. This agrees closely with the known mass of
\qty{0.050}{kg}, which supports $F = \frac{mv^2}{r}$.

% ---- Q13
\ansq
\begin{crit}
  \critrow{Compares velocity, acceleration AND net force for both motions, giving similarities and
    differences}{4}
  \critrow{Compares most of these quantities correctly}{3}
  \critrow{Describes some features of each motion}{2}
  \critrow{Gives one relevant point}{1}
\end{crit}
\sample
\begin{itemize}
  \item \textbf{Velocity:} in both, the direction of the velocity changes. In uniform circular motion
        the speed is constant; in projectile motion the speed changes (the horizontal component is
        constant but the vertical component changes).
  \item \textbf{Acceleration:} both accelerate. In circular motion the acceleration ($\frac{v^2}{r}$)
        is constant in size but always points to the centre, so its direction changes. A projectile's
        acceleration is $g$ downwards: constant in size and direction.
  \item \textbf{Net force:} in circular motion the net force points to the centre, at \ang{90} to the
        velocity, so it does no work. For a projectile the net force is its weight (constant), which
        is not always at \ang{90} to the velocity, so it does work and changes the speed.
\end{itemize}

% =====================================================================
\anstopic{Topic 3 - Motion in Gravitational Fields}

% ---- Q14
\ansq
\anspart{a}
\begin{crit}
  \critrow{Correctly calculates the force}{2}
  \critrow{Correct substitution into $F = \frac{GMm}{r^2}$}{1}
\end{crit}
\sample
\[ F = \frac{GMm}{r^2} = \frac{(\num{6.67e-11})(\num{6.0e24})(\num{7.35e22})}{(\num{3.84e8})^2} = \qty{2.0e20}{N} \]
\anspart{b}
\begin{crit}
  \critrow{Correctly calculates the orbital period in days}{2}
  \critrow{Uses $\frac{r^3}{T^2} = \frac{GM}{4\pi^2}$ (or calculates the orbital speed)}{1}
\end{crit}
\sample
\[ T = 2\pi\sqrt{\frac{r^3}{GM}} = 2\pi\sqrt{\frac{(\num{3.84e8})^3}{(\num{6.67e-11})(\num{6.0e24})}}
     = \qty{2.36e6}{s} = 27.4 \text{ days} \]
(or $v = \sqrt{\frac{GM}{r}} = \qty{1.02e3}{m.s^{-1}}$, then $T = \frac{2\pi r}{v}$.)

% ---- Q15
\ansq
\anspart{a}
\begin{crit}
  \critrow{Correctly calculates the altitude}{2}
  \critrow{Calculates the distance from Earth's centre}{1}
\end{crit}
\sample
$r = \sqrt{\dfrac{GM}{g}} = \sqrt{\dfrac{(\num{6.67e-11})(\num{6.0e24})}{4.9}} = \qty{9.04e6}{m}$.
Altitude $= \num{9.04e6} - \num{6.371e6} = \qty{2.7e6}{m}$. (Using $r = \sqrt{2}\,r_\mathrm{E}$, since
$g$ halves when $r^2$ doubles, gives \qty{2.6e6}{m}; also accepted.)
\anspart{b}
\begin{crit}
  \critrow{Draws a decreasing curve with an inverse-square shape AND shows correct values (about 9.8,
    2.5 and \qty{0.6}{N.kg^{-1}})}{2}
  \critrow{Draws a decreasing curve OR shows correct values}{1}
\end{crit}
\sample
\begin{center}
  \begin{tikzpicture}
    \draw[axisline] (0,0) -- (8.2,0) node[right] {$r$};
    \draw[axisline] (0,0) -- (0,5.6) node[above] {$g$ (\si{N.kg^{-1}})};
    \foreach \k/\l in {1/r_\mathrm{E}, 2/2r_\mathrm{E}, 3/3r_\mathrm{E}, 4/4r_\mathrm{E}}
      {\draw (\k*1.9,0.08) -- (\k*1.9,-0.08) node[below, font=\small] {$\l$};}
    \foreach \y/\l in {1/2, 2/4, 3/6, 4/8, 5/10}
      {\draw (0.08,\y*1.0) -- (-0.08,\y*1.0) node[left, font=\small] {\l};}
    \draw[line width=1.1pt, domain=1:4, samples=60, smooth] plot ({1.9*\x}, {4.9/(\x*\x)});
    \fill (1.9,4.9) circle (2pt) node[above right, font=\small] {9.8};
    \fill (3.8,1.225) circle (2pt) node[above right, font=\small] {2.45};
    \fill (7.6,0.306) circle (2pt) node[above, font=\small] {0.61};
  \end{tikzpicture}
\end{center}

% ---- Q16
\ansq
\anspart{a}
\begin{crit}
  \critrow{Calculates $\frac{r^3}{T^2}$ (or $\frac{T^2}{r^3}$) for at least three moons, or uses a ratio
    method for the whole data set, AND states that it is (nearly) constant, as Kepler's third law
    predicts}{2}
  \critrow{Calculates $\frac{r^3}{T^2}$ for one or two moons OR states Kepler's third law}{1}
\end{crit}
\sample
Kepler's third law says $\frac{r^3}{T^2}$ is the same for all moons of one planet.
Io: $\dfrac{(\num{4.22e8})^3}{(\num{1.53e5})^2} = \qty{3.21e15}{m^3.s^{-2}}$. Europa: \num{3.21e15}.
Ganymede: \num{3.21e15}. Callisto: \num{3.20e15}. The four values agree to within about 0.2\%, so
the data are consistent with the law.
\anspart{b}
\begin{crit}
  \critrow{Correctly calculates the mass of Jupiter}{2}
  \critrow{Uses $\frac{r^3}{T^2} = \frac{GM}{4\pi^2}$ with a correct value}{1}
\end{crit}
\sample
$M = \dfrac{4\pi^2}{G} \times \dfrac{r^3}{T^2} = \dfrac{4\pi^2 \times \num{3.21e15}}{\num{6.67e-11}}
= \qty{1.90e27}{kg}$.
\anspart{c}
\begin{crit}
  \critrow{Explains that the moon's mass cancels, so $\frac{r^3}{T^2}$ depends only on the mass of
    Jupiter}{1}
\end{crit}
\sample
The gravitational force and the centripetal force are both proportional to the moon's mass, so it
cancels: $\frac{r^3}{T^2} = \frac{GM}{4\pi^2}$ depends only on Jupiter's mass $M$.

% ---- Q17
\ansq
\anspart{a}
\begin{crit}
  \critrow{Correctly calculates the orbital speed}{2}
  \critrow{Calculates the orbital radius OR uses $v = \sqrt{\frac{GM}{r}}$}{1}
\end{crit}
\sample
$r = \num{6.371e6} + \num{6.0e5} = \qty{6.971e6}{m}$.
$v = \sqrt{\dfrac{GM}{r}} = \sqrt{\dfrac{(\num{6.67e-11})(\num{6.0e24})}{\num{6.971e6}}} = \qty{7.58e3}{m.s^{-1}}$.
\anspart{b}
\begin{crit}
  \critrow{Correctly calculates the total energy (negative)}{2}
  \critrow{Uses $E = -\frac{GMm}{2r}$ OR calculates $K$ and $U$}{1}
\end{crit}
\sample
\[ E = -\frac{GMm}{2r} = -\frac{(\num{6.67e-11})(\num{6.0e24})(1200)}{2 \times \num{6.971e6}} = \qty{-3.44e10}{J} \]
(or $K = \frac{1}{2}mv^2 = \qty{3.44e10}{J}$ and $U = -\frac{GMm}{r} = \qty{-6.89e10}{J}$, so
$E = K + U = \qty{-3.44e10}{J}$).
\anspart{c}
\begin{crit}
  \critrow{Explains that $U$ is negative (zero at infinity) and larger in size than $K$, so the satellite
    is bound to Earth}{1}
\end{crit}
\sample
$U$ is zero at infinity and negative everywhere else. In orbit $|U| = 2K$, so $E = K + U$ is
negative. This means the satellite is bound to Earth: energy must be added to it before it can
escape.

% ---- Q18
\ansq
\anspart{a}
\begin{crit}
  \critrow{Correctly calculates the speed at aphelion}{2}
  \critrow{Writes $r_\mathrm{p}v_\mathrm{p} = r_\mathrm{a}v_\mathrm{a}$ with correct values}{1}
\end{crit}
\sample
$r_\mathrm{p}v_\mathrm{p} = r_\mathrm{a}v_\mathrm{a}$, so
$v_\mathrm{a} = \dfrac{(\num{1.0e11})(30)}{\num{1.5e11}} = \qty{20}{km.s^{-1}}$.
\anspart{b}
\begin{crit}
  \critrow{Explains that the total energy is constant AND that $U$ increases (becomes less negative)
    as $r$ increases, so $K$ and the speed decrease}{2}
  \critrow{Provides some relevant information}{1}
\end{crit}
\sample
Only gravity acts, so the total energy $K + U$ stays constant. As the planet moves away from the
star, $r$ increases, so $U = -\frac{GMm}{r}$ becomes less negative (increases): the planet does work
against gravity. For the total to stay the same, $K$ must decrease, so the planet slows down.

% ---- Q19
\ansq
\anspart{a}
\begin{crit}
  \critrow{Correctly calculates the escape velocity}{1}
\end{crit}
\sample
$v_\mathrm{esc} = \sqrt{\dfrac{2GM}{r}} = \sqrt{\dfrac{2(\num{6.67e-11})(\num{6.0e24})}{\num{6.371e6}}} =
\qty{1.12e4}{m.s^{-1}}$.
\anspart{b}
\begin{crit}
  \critrow{Shows (with the energy equation) that both $K$ and $U$ are proportional to the rocket's mass,
    so the mass cancels}{2}
  \critrow{Provides some relevant information}{1}
\end{crit}
\sample
To escape, $\frac{1}{2}mv^2 - \frac{GMm}{r} \ge 0$. Both the kinetic energy and the potential energy are
proportional to the rocket's mass $m$, so $m$ cancels: $v_\mathrm{esc} = \sqrt{\frac{2GM}{r}}$. A
heavier rocket needs more \emph{energy} to escape, but the same \emph{speed}.
\anspart{c}
\begin{crit}
  \critrow{Explains that the satellite only needs the speed for a circular orbit
    ($\sqrt{\frac{GM}{r}}$, which is $\frac{1}{\sqrt{2}}$ of the escape velocity) because it stays bound
    to Earth and does not travel to infinity}{1}
\end{crit}
\sample
Escape velocity is the speed needed to travel to infinity and never come back. A satellite in orbit
stays bound to Earth (its total energy is negative). It only needs the speed at which gravity
provides exactly the centripetal force, $v = \sqrt{\frac{GM}{r}}$, which is $\frac{1}{\sqrt{2}}$ of the
escape velocity at the same distance.

% ---- Q20
\ansq
\anspart{a}
\begin{crit}
  \critrow{Correctly calculates the orbital radius}{3}
  \critrow{Substitutes correctly into $r^3 = \frac{GMT^2}{4\pi^2}$ with $T$ in seconds}{2}
  \critrow{Converts the period to seconds OR states the correct relationship}{1}
\end{crit}
\sample
$T = 24.6 \times 3600 = \qty{88560}{s}$.
\[ r = \sqrt[3]{\frac{GMT^2}{4\pi^2}} = \sqrt[3]{\frac{(\num{6.67e-11})(\num{6.42e23})(88\,560)^2}{4\pi^2}}
   = \qty{2.04e7}{m} \]
\anspart{b}
\begin{crit}
  \critrow{Correctly calculates the altitude}{1}
\end{crit}
\sample
Altitude $= \num{2.04e7} - \num{3.39e6} = \qty{1.70e7}{m}$.
\anspart{c}
\begin{crit}
  \critrow{Explains that the satellite must turn with Mars (same period and direction) AND that only an
    equatorial orbit keeps it above one point}{2}
  \critrow{Provides some relevant information}{1}
\end{crit}
\sample
To stay above one point, the satellite must go around once in the same time as Mars rotates, and in
the same direction, so it keeps pace with the surface. Every orbit is centred on the centre of Mars.
An orbit tilted to the equator would carry the satellite north and south of the equator during each
orbit, so only an orbit directly above the equator keeps it above the same point.

\end{document}
